% =====================================================================
% ROZDZIAŁ 1 — Wstęp
% =====================================================================
\chapter{Wstęp}
\label{rozd:wstep}

\section{Problem i motywacja}
\label{sec:w-problem}

W poniedziałek rano przejmowane jest konto pocztowe pracownika. Zanim
doba się skończy, jego kontakty w dwóch działach i w firmie partnerskiej
otrzymują krótką, wiarygodną wiadomość z linkiem --- każda od kolegi, z
którym korespondują co tydzień; część z nich klika, a godzinę później
ich własne kontakty otrzymują tę samą wiadomość. Żaden pojedynczy
element żadnej z tych wiadomości nie zdradza oszustwa
\citep{butavicius2022why,sturman2025security}. To \emph{lateral phishing}:
kampania wysłana z prawdziwego, wewnętrznego konta przejętego przez
atakującego; w skali systemowej scharakteryzowali go \citet{ho2019detecting}
jako ataki docierające do wielu odbiorców wewnątrz i wokół organizacji.
Raporty FBI IC3 szacują straty spowodowane przez Business Email Compromise
(BEC), którego lateral phishing jest odmianą „od środka", na około
2{,}9 miliarda dolarów w~2023 roku \citep{fbi2024ic3}.

W drugim scenariuszu, lustrzanym wobec pierwszego, na skrzynkę dyrektora
finansowego trafia prośba o pilny przelew od „prezesa". Treść jest
dopracowana, pretekst wiarygodny, a nadawca --- zewnętrzny --- respektuje
wszystko, czego adresat oczekuje: właściwe nazwisko, właściwy projekt,
właściwy ton \citep{dewan2014analyzing,brown2025web}; koszt takiej
personalizacji językowe modele obniżają z każdym rokiem
\citep{bethany2024evaluating}. To \emph{impersonacja zewnętrzna}, czyli
spear phishing w reżimie BEC.

Detekcja phishingu bywa przedstawiana jako problem niemal rozwiązany.
Nowoczesne filtry, oparte na treści i reputacji, klasyfikują kampanie
masowe z dokładnością $\sim$98\% na typowych benchmarkach
\citep{birthriya2025detection,tan2020graph}. Liczba ta jest jednak
myląca: dokładność na benchmarkach jest zdominowana przez ogon masowy,
podczas gdy rzadki ogon ukierunkowany --- spear phishing i BEC ---
napędza większość strat finansowych \citep{cidon2019high,ho2019detecting},
jednocześnie z definicji omijając metody oparte na treści. Lateral phishing
jest najtrudniejszym przypadkiem tego ogona: przejęte konto przechodzi
każdą kontrolę tożsamości, reputacji czy relacji, bo wszystkie one są
prawdziwe. W obu scenariuszach oszustwo nie leży w tym, \emph{co} jest
napisane, lecz w tym, \emph{kto} pisze \emph{kiedy} --- i co dzieje się
potem.

\section{Teza rozprawy}
\label{sec:w-teza}

Rozprawa broni następującej tezy, sformułowanej w dwóch częściach.

\begin{description}
\item[T1 (sygnał).] Phishing ukierunkowany --- zarówno w wariancie lateral
(przejęcie konta wewnętrznego), jak i w wariancie impersonacji
zewnętrznej --- jest mierzalnym naruszeniem \emph{grafów organizacji}:
grafu komunikacji (kto koresponduje z kim, w jakim rytmie) oraz grafu
wiedzy domenowej (wielowarstwowa struktura relacji: kontakty,
hierarchia, kompetencje, publiczny ślad OSINT, rytm czasowy). Sygnał
detekcyjny dla ogona ukierunkowanego mieszka w dynamice czasowej i
strukturalnej tych grafów, a nie w treści pojedynczej wiadomości:
(a)~odbiorcza fala kaskady jest widoczna dla detektora temporalnego
i degraduje łagodniej niż detekcja wolumetryczna pod atakiem skrytym,
który rozkłada wysyłki w czasie; (b)~niespójność międzywarstwowa
multipleksu wiedzy domenowej oddziela impersonację zewnętrzną od
legalnych nowych nadawców --- tam, gdzie detektor oparty wyłącznie na
kontaktach zawodzi całkowicie.
\item[T2 (ewaluacja).] Empiryczny wynik detektora strumieniowego --- AUC,
czułość w punkcie pracy --- mówi cokolwiek o zachowaniu w warunkach
produkcyjnych tylko wtedy, gdy spełnione są cztery warunki
przyczynowości: (i)~cechy są zdarzeniowo-przyczynowe (nie widzą
przyszłości), (ii)~podział ewaluacyjny jest rozłączny po ofiarach,
(iii)~brzegowa sygnatura czasowa iniekcji jest zmierzona i usunięta,
(iv)~fabrykacja śladów jest losowana z rozkładu śladów prawdziwych.
Naruszenie każdego z tych warunków --- w literaturze i we wcześniejszych
wersjach własnych eksperymentów --- powodowało zawyżenie raportowanych liczb; rozprawa
pokazuje to formalnie (twierdzenia z konstrukcyjnymi kontrprzykładami,
rozdz.~\ref{rozd:formalny}) i empirycznie (pułapka wycieku rytmu:
pozorna czułość 0{,}73 spada do 0{,}00 pod uczciwym podziałem
czasowym).
\end{description}

Obie części łączy wspólna metoda: każda liczba w rozprawie pochodzi
z paczki reprodukcyjnej (\texttt{code/results/*.csv}), każda przewaga
detektora jest poddana ablacyjnej kontroli (tasowanie czasu, tasowanie
warstw, sweep fabrykacji), a każde twierdzenie operacyjne jest
wypowiedziane z jawną arytmetyką bazy (prewalencją): przy realnej
częstości ataku rzędu $1{:}10^4$ klasa detektorów jest narzędziem
\emph{rankingowym} dla analityków SOC, nie autonomiczną obroną.

\section{Pytania badawcze}
\label{sec:w-rq}

Teza jest testowalna przez cztery pytania badawcze --- po dwa na każdą
część tezy i po jednym na każdy z artykułów stanowiących podstawę
rozdziałów empirycznych.

\begin{description}
\item[RQ1 (sygnał kaskadowy).] Czy przejęte konto, zgodne z każdą
statyczną relacją grafu komunikacji, da się wykryć z dynamiki kaskady ---
i czy sygnał przeżywa atakującego skrytego, który rozkłada wysyłki
w czasie? Podteza: detektor temporalny czytający odbiorczą falę bije
detekcję wolumetryczną w niskim FPR i degraduje pod skrytym atakiem
mniej; przewaga znika pod kontrolą tasowania czasu i przenosi się
kierunkowo na realne topologie (rozdz.~\ref{rozd:kaskada}).
\item[RQ2 (granice sygnału kaskadowego).] Jakie są operacyjne granice
tego sygnału? Podteza: sygnał wymaga fan-outu $\ge2$ (atak chirurgiczny
na jednego odbiorcę jest niewidoczny dla wszystkich detektorów),
degraduje łagodnie pod szumem etykiet i skalą, a przy realnej
prewalencji detektory służą jako instrument triażu (rozdz.~\ref{rozd:kaskada}).
\item[RQ3 (sygnał międzywarstwowy).] Czy impersonacja zewnętrzna
oddziela się od legalnych nowych nadawców przez \emph{niespójność
międzywarstwową} grafu wiedzy domenowej --- i czy zysk pochodzi ze
struktury warstw, a nie z pojemności cech? Podteza: fuzja warstw z rytmem
podnosi AUC i likwiduje fałszywe alarmy na legalnych nowych nadawcach,
kontrola tasowania sprowadza wynik do losowego, a realny multipleks
Enronu zachowuje operacyjną czułość pod uczciwym podziałem czasowym
(rozdz.~\ref{rozd:multipleks}).
\item[RQ4 (narażenie i granice).] Co graf wiedzy domenowej mówi o tym,
\emph{kto} jest najbardziej narażony, i co się dzieje z tą klasą
detektorów, gdy ich ewaluacja przecieka? Podteza: kanał ataku, nie ranga
ofiar w organizacji, napędza narażenie (zimna impersonacja zostawia
strukturową podłogę $\mathrm{Susc}\approx0{,}96$ na każdym stanowisku),
a wyciek rytmu powoduje zawyżenie czujności detektorów behawioralnych o rzędy
wielkości (rozdz.~\ref{rozd:multipleks}).
\end{description}

\section{Wkład rozprawy}
\label{sec:w-wklad}

Wkład naukowy rozprawy składa się z czterech elementów, mapowanych na
rozdziały.

\begin{enumerate}
\item \textbf{Rozdział formalny} (rozdz.~\ref{rozd:formalny}):
aksjomatyzacja warunków ewaluacji strumieniowej --- klasa cech i
detektorów zdarzeniowo-przyczynowych, cztery mechanizmy wycieku
(look-ahead, cechy definiujące etykietę, sygnatura iniekcji, ślepota
wewnątrzkubełkowa) jako twierdzenia z konstrukcyjnymi kontrprzykładami,
dekompozycja wyniku na składową czysto czasową i sprzęgającą czas
z grafem, wynik o nierozróżnialności śladów przy realistycznej
fabrykacji oraz arytmetyka bazy.
\item \textbf{Formalizacja i pomiar naruszenia grafu komunikacji}
(rozdz.~\ref{rozd:kaskada}, na bazie artykułu A, \emph{Journal of
Network and Computer Applications}): lateral phishing jako temporalne
naruszenie grafu komunikacji organizacji, z modelem zagrożeń obejmującym
kaskadę naiwną i atakującego skrytego; odbiorcza fala kaskady jako
niezmiennik detekcyjny --- hierarchia detektorów, kontrola tasowania
czasu, wynik główny (detekcja wolumetryczna załamuje się poniżej losowego
podstawienia pod atakiem skrytym, sygnał odbiorczy degraduje mniej
i trzyma się powyżej losowego); replikacja mechanizmu na trzech realnych
grafach komunikacji (Enron, \texttt{email-Eu-core}, \texttt{CollegeMsg}).
\item \textbf{Formalizacja i pomiar naruszenia grafu wiedzy domenowej}
(rozdz.~\ref{rozd:multipleks}, na bazie
artykułu B, \emph{IEEE}): graf wiedzy domenowej jako multipleks,
niespójność międzywarstwowa jako niezmiennik detekcyjny dla
impersonacji zewnętrznej; odkrycie metodyczne --- \emph{pułapka wycieku
rytmu} --- wraz z detektorowo-relatywną miarą narażenia
(susceptibility) i granicą graf--treść.
\item \textbf{Metodologia ewaluacji typu leak-aware}: protokół
czterech zabezpieczeń (dopasowany projekt iniekcji, kontrola
tasowania, podział indukcyjny po organizacjach, audyt szumu etykiet)
zoperacjonalizowany w rozdziale formalnym i zastosowany w obu
rozdziałach empirycznych; operacyjna definicja tego, co znaczy „detektor
tej klasy działa", wraz z jawnym wyznaczeniem granic (punkt ślepy
fan-out 1, szum etykiet, skala, próg bazowy).
\end{enumerate}

Modele uczone użyte w rozprawie --- temporal graph network z pamięcią
węzłów i gradient boosting --- to standardowe narzędzia; wkładem jest
ujęcie phishingu ukierunkowanego jako mierzalnego naruszenia grafów
organizacji oraz dowód, że sygnał odbiorczy zależy od porządku
czasowego, przenosi się kierunkowo na realne grafy i degraduje pod
skrytością mniej niż detekcja wolumetryczna.

\section{Powiązanie z artykułami}
\label{sec:w-artykuly}

Rozprawa jest zbudowana wokół dwóch powiązanych artykułów, których
pełne teksty (po angielsku) stanowią podstawę rozdziałów empirycznych:

\begin{itemize}
\item \textbf{Artykuł A} (journal, \emph{Journal of Network and
Computer Applications}): kaskadowy lateral phishing, temporal GNN
przeciwko detektorom wolumetrycznym, kontrola time-shuffle, walidacja
Enron + SNAP --- materiał rozdziału~\ref{rozd:kaskada};
\item \textbf{Artykuł B} (konferencja IEEE): multiplex grafu wiedzy
domenowej, niespójność międzywarstwowa, pułapka wycieku rytmu na
Enronie, narażenie detektorowo-relatywne --- materiał
rozdziału~\ref{rozd:multipleks};
\item \textbf{Artykuł C} (konferencja IEEE, roboczy): graf wiedzy jako
narzędzie operacyjne (note-grafy, CTI KG, GraphRAG) i jako cel ataku ---
realistyczna fabrykacja śladów OSINT, trucie korpusu GraphRAG-lite
z bramką spójności, fabrykacja trójek na publicznym FB15k-237 ---
materiał sekcji~\ref{sec:m-kg} rozdziału~\ref{rozd:multipleks}.
\end{itemize}

Każde twierdzenie rozdziału formalnego mapuje się na konkretny
eksperyment z paczki reprodukcyjnej (tab.~\ref{tab:mapping}
w rozdz.~\ref{rozd:formalny}); każda liczba w rozdziałach empirycznych
pochodzi wprost z plików \texttt{code/results/*.csv}. Dwie wcześniejsze
usterki wycieku (momentum liczone z przyszłych kubełków, liczniki
wolumetryczne liczone na całym strumieniu, cechy stopni GNN ze splitu
testowego) zostały usunięte przed przeliczeniem \emph{wszystkich}
wyników --- zgodnie z regułą: każda cecha strumieniowa liczona
wyłącznie z przeszłości, \emph{zgodnie z pamięcią modelu temporalnego}
(rozdz.~\ref{rozd:formalny}, def.~\ref{def:causal-feature}).

\section{Struktura rozprawy}
\label{sec:w-struktura}

Rozdział~\ref{rozd:przeglad} przegląda literaturę: detekcję phishingu
masowego i ukierunkowanego, detekcję anomalii w poczcie
behawioralnej, sieci multiplex i grafowe modele uczone oraz rygor
ewaluacji w bezpieczeństwie ML. Rozdział~\ref{rozd:formalny} buduje
ramę formalną: porządek totalny zdarzeń, przyczynowość zdarzeniową,
cztery mechanizmy wycieku, sondę czasową, kontrolę tasowania,
realistyczną fabrykację i arytmetykę bazy --- z mapowaniem na
eksperymenty. Rozdział~\ref{rozd:kaskada} (na bazie artykułu A)
rozwija lateral phishing jako naruszenie grafu komunikacji w czasie:
model, detektory, wyniki na grafie studyjnym, atakujący skryty,
granice modelu zagrożeń, walidacja na realnych grafach i rygor
ewaluacji. Rozdział~\ref{rozd:multipleks} (na bazie artykułu B)
rozwija impersonację zewnętrzną jako naruszenie grafu wiedzy
domenowej: multipleks, niespójność międzywarstwową, pułapkę wycieku
rytmu na realnym Enronie, narażenie detektorowo-relatywne i granicę
graf--treść. Rozdział~\ref{rozd:synteza} syntetyzuje wyniki obu
rozdziałów empirycznych w świetle twierdzeń formalnych i wypowiada
implikacje operacyjne dla SOC. Rozdział~\ref{rozd:zakonczenie}
podsumowuje odpowiedzi na pytania badawcze, ograniczenia i kierunki
przyszłej pracy.
